\section{Cash exercise across one unrevealed meeting}\label{app:crossmeeting}
This formula uses the pure-jump, zero-initial-curve model and supplies an explicit stopping rule, rather than only an equality of values across allocations. Let
$T_{N-1}\leq A<T_N<S\leq a<b$, set $T_*=T_N$, $\ell=T_*-A$, $L=S-T_*$, and $\bar t=(a+b)/2$. Write $V_-=\sum_{i<N}v_i$, $v_*=v_N$, and $V=V_-+v_*$. Define
\[
f_-(x)=\frac{\exp\{\delta[x+V_-(\bar t-A)+v_*(b-T_*)]\}-1}{\delta},
\qquad
f_+(y)=\frac{\exp\{\delta[y+V(\bar t-T_*)]\}-1}{\delta}.
\]
The futures rate equals $f_-(r_A)$ before $T_*$ and $f_+(r_{T_*})$ afterwards. Extend $\mathcal D_{V,L}$ from Section~\ref{sec:aggregation} to any variance $q\geq0$ and duration $d\geq0$ by
\[
\mathcal D_{q,d}(x)=\max_{0\leq u\leq d}e^{-xu-qu^2/2}.
\]
For a fixed rate strike $k$, put
\begin{align*}
p_*(x)&=(f_-(x)-k)^+\mathcal D_{V_-,\ell}(x),\\
c_*(x)&=e^{-x\ell-V_-\ell^2/2}
  \E_{Z\sim N(0,v_*)}\left[
       (f_+(x+V_-\ell+Z)-k)^+
                  \mathcal D_{V,L}(x+V_-\ell+Z)\right].
\end{align*}
The American cash value is
\begin{equation}\label{eq:crosscash}
U=\E_{Y\sim N(0,V_-)}[\max\{p_*(Y),c_*(Y)\}].
\end{equation}
This is a one-dimensional Gaussian integral containing one further Gaussian expectation. It depends on the earlier allocation only through $V_-$, with $v_*$ separately retained.

To prove the formula, condition on $r_A=x$. Before $T_*$, the filtration contains no new information and $r_{A+u}=x+V_-u$. Every discounted exercise payoff is bounded by $p_*(x)$. Waiting for the announcement yields the continuation value $c_*(x)$ after optimally choosing the cash-discount timing on $[T_*,S]$. The futures martingale identity gives
$\E[f_+(x+V_-\ell+Z)]=f_-(x)$, hence
\[
c_*(x)\geq e^{-x\ell-V_-\ell^2/2}(f_-(x)-k)^+.
\]
If the maximizer defining $p_*(x)$ lies at $\ell$, this inequality implies $c_*(x)\geq p_*(x)$. Thus whenever $p_*(x)>c_*(x)$, its maximizing exercise time is strictly before the unrevealed meeting and is admissible. Exercise then; otherwise wait, with ties assigned to waiting, and optimize after revelation. This rule is a stopping time and attains the conditional maximum. The $A$-discount tilt gives $r_A\sim N(0,V_-)$ and proves \eqref{eq:crosscash}. Gaussian exponential moments and $\mathcal D_{q,d}(x)\leq e^{d|x|}$ give finiteness, including zero variances.

Bond-implied discrete compounding gives the same formulas when the fixing partition contains every meeting inside accrual; in this case all meetings precede accrual, so the product telescopes automatically. This supplies an exact example in which the exercise window crosses an information release without resolving the allocation among earlier meetings.
